\section{Introduction}
\label{sec:intro}

Two research communities have arrived, independently, at the same
weight alphabet. Ternary large language models --- BitNet-class
systems shipped in $\{-1, 0, +1\}$ quantizations with integer
runtimes \citep{bitnet158,bitnnetcpp} --- put a three-symbol weight
world into production. Spiking neural networks compute natively in
that same integer regime, and their local, online plasticity rules
(STDP, Hebbian learning) update weights without backpropagation,
gradients, or loss-guided search. The literatures do not cite each
other, and --- as of this paper's surveyed date, with the field's own
surveys marking the gap open \citep{surveylargesnn,rajendran2023future} ---
no work has joined them: no published system applies local plasticity
to a shipped production LLM's weights.

This paper is the record of joining them, with the honesty that
turned out to require. We present \system{}, a
\texttt{no\_std} \texttt{i16} fixed-point spiking substrate whose
format bridge speaks the shipped ternary formats byte-exactly
(Sec.~\ref{sec:substrate}--\ref{sec:bridge}), and the closed loop it
enables (Sec.~\ref{sec:loop}): import a real 4B Q2\_0 LLM's weight
tensor into the spiking substrate, drive the substrate with
\emph{the model's own activations} (Sec.~\ref{sec:invivo}), let local
plasticity adapt the imported weights, re-export through the bridge
byte-exactly into a copy of the shipped file, and hand that file to a
pinned foreign runtime --- a llama.cpp fork --- which executes it.
Every step is deterministic and byte-audited; a control export
carries unadapted trits through the identical machinery and comes out
byte-identical to the original.

The loop closes. The adapted file measurably shifts the foreign
runtime's logits (60/60 generation steps at every judged footprint),
and the in-vivo-adapted file changes greedy continuations on 2 of 5
frozen prompts. The question that decides whether any of that is a
\emph{capability} is: do those changes carry the adaptation's
content, or would equal-dose random perturbation do the same? We
answered it with a pre-registered null ladder --- margin census,
dose-matched random families, a value-flip family, sibling prompts, a
report-only stress arm, three sealed rules --- committed in full
before any null table was read (Sec.~\ref{sec:adjudication}). The
verdict, by rule: the changes failed every pre-registered criterion.
Eight of ten equal-dose random perturbations land the in-vivo run's
exact destination. The term of record is \unattrib{}.

We report that negative as a first-class result for three reasons.
First, it was \emph{our} false positive to catch: the project's own
ledger, pre-adjudication, recorded the continuation change as a
capability demonstration (quoted verbatim in
Sec.~\ref{sec:averted}); only the pre-committed rule reversed the
authors' first reading. Flip-count evidence in this regime publishes
false positives. Second, the negative is structurally informative:
the model has generic destination basins that equal-dose noise
reaches at high probability, its knife-edges are fragile to
perturbation magnitude, and --- by a theorem this paper proves from
the terminal diff's composition --- content and placement are
formally entangled at this composition, so the natural control rung
does not exist (Sec.~\ref{sec:discussion}). Third, what survives is
exactly the infrastructure claim the title makes: a spiking substrate
that tracks a shipped quantized LLM's arrangement, adapts under its
own activations, and round-trips its weights through the production
format into foreign tooling --- first as surveyed, with the survey's
boundaries stated and the pre-submission re-verification standing as
a gate (Sec.~\ref{sec:related}).

\paragraph{Contributions.}
\begin{itemize}
\item \textbf{The substrate and bridge} (Sec.~\ref{sec:substrate},
  \ref{sec:bridge}): a published, \texttt{no\_std}, integer-only SNN
  library with byte-exact two-way ternary format bridges, and the
  falsification-driven discovery of a structural transmission bug
  whose fix reversed the substrate's plasticity mechanism ---
  a case study in pre-registered debugging inside infrastructure.
\item \textbf{The loop} (Sec.~\ref{sec:loop}--\ref{sec:invivo}): the
  L1--L4 chain closed on a shipped 4B Q2\_0 LLM --- local plasticity
  under the model's own drive, byte-audited re-export, foreign
  execution, control identity --- with the substrate's
  arrangement-tracking result under both synthetic and in-vivo drives
  (n=1, caveats in-text).
\item \textbf{The adjudication} (Sec.~\ref{sec:adjudication}): a
  pre-registered null ladder that refutes content-attribution at this
  composition and reports \unattrib{} --- including the two rungs we
  believe are novel: a composition-specific null-feasibility theorem,
  and destination-basin analysis, which together caught the false
  positive the authors' own first reading had become.
\end{itemize}

The remainder of the paper: Discussion of what would have counted and
why the negative informs (Sec.~\ref{sec:discussion}); related work and
the survey matrix (Sec.~\ref{sec:related}); limitations, pre-registered
(Sec.~\ref{sec:limitations}); destination transcripts and full
reproducibility material in the appendices.
